\section{Literature Review and Research Positioning}
\label{sec:literature-review}
The relevant literature connects numerical resource optimization, stateful reconfiguration, and agentic operational diagnosis. These strands already provide many of the mechanisms that an adaptive streaming controller would require. The purpose of this review is to distinguish those inherited capabilities from the proposed evidence-acquisition question, rather than to treat their combination as automatically novel. Tables~\ref{tab:related-work} and~\ref{tab:additional-related-work} position thirty published papers, including fifteen additions to the initial review. Comparisons below concern the inspected mechanisms and evaluation scope; they are not controlled reruns of the published systems.

\begin{table*}[t]
\caption{Positioning of the initial fifteen reviewed papers. The final column identifies how each source informs EcoStream; it does not claim that the source lacks every proposed feature. Scientific agents and power orchestration are supporting references rather than directly interchangeable streaming baselines.}
\label{tab:related-work}
\centering
\footnotesize
\setlength{\tabcolsep}{5pt}
\renewcommand{\arraystretch}{1.13}
\begin{tabularx}{\textwidth}{@{}p{0.19\textwidth} >{\raggedright\arraybackslash}X >{\raggedright\arraybackslash}X@{}}
\toprule
\textbf{Paper / system} & \textbf{Established mechanism} & \textbf{Role in the present study} \\
\midrule
ORACL~\cite{bai2026oracl} & LLM diagnosis, counterfactual allocation, and guarded execution & Closest joint diagnosis--allocation overlap \\
RCAgent~\cite{wang2024rcagent} & Autonomous diagnostic tools, observation management, and finalization & Conventional agentic diagnosis and stopping comparator \\
Demeter~\cite{geldenhuys2024demeter} & Multi-configuration optimization and selective profiling & Numerical optimization and information-cost accounting \\
STREAMLINE~\cite{pedratscher2025streamline} & Ensemble-wide auto-tuning, scheduling, and state-preserving updates & Recent holistic adaptation precedent \\
CAOM~\cite{tan2025caom} & Bottleneck detection and cost-aware migration timing & Transition timing and backlog-aware decisions \\
MUDAP/RASK~\cite{sedlak2026mudap} & Structural regression and constrained multi-dimensional adaptation & Sample-efficient numerical alternative; action-space caution \\
DRRS~\cite{qing2025drrs} & Fine-grained stateful scaling execution & Realistic transition-cost and state-size evaluation \\
Ra-Stream~\cite{fan2026rastream} & Communication-aware partitioning, scheduling, and scaling & Non-capacity bottlenecks and topology context \\
SAP-DPO~\cite{ren2026sapdpo} & Dynamic partitioning with adaptive update frequency & Skew-aware interpretation and monitoring precedent \\
ESSA/TTCLB-EA~\cite{li2025energyflink} & SLA-based energy-aware scheduling and coordinated load balancing & Efficiency motivation; resource-versus-energy distinction \\
Daedalus~\cite{pfister2024daedalus} & Online capacity models and proactive horizontal scaling & Low-overhead controller and service/resource trade-offs \\
Phoebe~\cite{geldenhuys2022phoebe} & Profiling and forecasting with recovery constraints & Recovery-aware attribution and backlog-clearance requirements \\
SciToolAgent~\cite{ding2025scitoolagent} & Knowledge-graph-assisted scientific tool selection & Bounded contracts and justified tool orchestration \\
BioMedAgent~\cite{bu2026biomedagent} & Interactive workflow execution and experience retrieval & Outcome feedback and separately evaluated adaptation \\
Grid-interactive AI~\cite{colangelo2026grid} & Measured power orchestration under flexibility-tier commitments & Physical measurement boundaries and workload constraints \\
\bottomrule
\end{tabularx}
\end{table*}

\begin{table*}[t]
\caption{Fifteen additional published papers and their comparative roles. Fourteen are conventional journal articles; Scarf appears in the journal-style Proceedings of the VLDB Endowment. Two entries are surveys, not implementation baselines. No addition is cited as an arXiv preprint.}
\label{tab:additional-related-work}
\centering
\footnotesize
\setlength{\tabcolsep}{5pt}
\renewcommand{\arraystretch}{1.13}
\begin{tabularx}{\textwidth}{@{}p{0.19\textwidth} >{\raggedright\arraybackslash}X >{\raggedright\arraybackslash}X@{}}
\toprule
\textbf{Paper / system} & \textbf{Established mechanism} & \textbf{Role in the present study} \\
\midrule
Scarf~\cite{liu2026scarf} & Multi-objective RL, knob selection, and transferable tuning & Strong recent Flink tuning comparator \\
MMP-SPS~\cite{wladdimiro2026mmpsps} & Multi-window prediction with adaptive model selection & Forecasting-only alternative; overhead and service trade-offs \\
Ns-Stream~\cite{wu2025nsstream} & Near-source grouping to reduce communication & Communication-versus-capacity diagnosis; engine boundary \\
Ms-Stream~\cite{wu2025msstream} & Hierarchical grouping, partitioning, and placement & Task/node stragglers and valid remediation actions \\
ProOSStream~\cite{lajili2026proosstream} & Offline optimization, federated RL, and forecasting & Distributed scheduling precedent, not single-host validation \\
Hierarchical autoscaling~\cite{russorusso2023hierarchical} & Two-layer BO/RL control of heterogeneous resources & Strong numerical control and learning-budget comparison \\
CEBFD~\cite{li2024cebfd} & Demand prediction and cost-aware Spark executor placement & Demand/placement separation; micro-batch scope \\
Local recovery~\cite{takdir2025localrecovery} & Partial snapshots and subset restoration in Flink & Actual recovery scope and replay-cost accounting \\
A-FP4S~\cite{xu2023afp4s} & Adaptive fragment-based parallel state reconstruction & Recovery-cost precedent; not an EcoStream invention \\
TAMO~\cite{zhang2025tamo} & Published tool-assisted multimodal LLM diagnosis scope & Close overlap; final method access still pending \\
LocaleXpert~\cite{zhong2026localexpert} & Modality interpretation and specialist diagnostic agents & Multimodal diagnosis control; final-version check pending \\
MA-RCA~\cite{fu2026marca} & Retrieval grounding and runtime hypothesis validation & Validation overhead and conventional agent comparison \\
CausalRCA~\cite{xin2023causalrca} & Graph-based fine-grained causal metric ranking & Non-LLM diagnosis comparator; action effects separate \\
Stream evolution survey~\cite{fragkoulis2024evolution} & State, ordering, fault tolerance, and elasticity overview & Systems background and semantic constraints \\
Runtime adaptation survey~\cite{cardellini2022adaptation} & Taxonomy of stream-processing adaptation strategies & Coverage framework, not an executable baseline \\
\bottomrule
\end{tabularx}
\end{table*}

Two surveys provide complementary foundations. Fragkoulis et al. connect stream-processing evolution to state, ordering, fault tolerance, elasticity, and reconfiguration~\cite{fragkoulis2024evolution}. Cardellini et al. organize runtime adaptation through a taxonomy of approaches and open challenges~\cite{cardellini2022adaptation}. They motivate treating adaptation as an established systems problem with several interacting constraints. As surveys, they inform terminology and coverage rather than constitute executable controller baselines.

\subsection{Numerical Autoscaling and Configuration Optimization}
Phoebe combines workload forecasting, parallel profiling, and quality-of-service models to tune streaming jobs under latency and recovery requirements~\cite{geldenhuys2022phoebe}. Its recovery discussion establishes an important constraint: capacity sufficient to match normal input can still be insufficient to clear the backlog after a failure. Daedalus instead builds processing-capacity models from runtime observations and uses forecasting for proactive horizontal scaling~\cite{pfister2024daedalus}. Its comparisons illustrate that reducing resources can involve a latency trade-off, rather than improving every service metric simultaneously.

Demeter expands the decision space to worker count, CPU, memory, processing slots, and checkpoint interval, combining forecasting with multi-objective Bayesian optimization~\cite{geldenhuys2024demeter}. It already selects additional configuration-profiling runs according to information gain and skips unnecessary profiling. Consequently, selective acquisition of information in general is not an unoccupied research area. The candidate distinction for EcoStream concerns operational diagnostic observations that separate competing bottleneck explanations, rather than additional executions intended to learn a configuration's performance.

More recent work broadens adaptation beyond isolated job parameters. STREAMLINE coordinates workload prediction, evolutionary configuration tuning, and operator scheduling across pipeline ensembles, with a state-preserving update mechanism described in its publisher abstract~\cite{pedratscher2025streamline}. MUDAP exposes service- and resource-level elasticity parameters, while its RASK agent uses structural knowledge, regression, and a constrained solver to coordinate competing services~\cite{sedlak2026mudap}. Some MUDAP actions alter service quality or model size. Such flexibility is valuable in its setting, but cannot be assumed equivalent to changing resources while preserving an exact streaming query. EcoStream must therefore compare controllers under the same permitted actions and semantics, including a strong low-overhead numerical alternative.

Scarf provides a particularly relevant recent Flink comparator. It combines workload-sensitivity-based knob selection, multi-objective reinforcement learning, and transferable graph-based models to tune throughput and resource usage~\cite{liu2026scarf}. Its publisher-deposited abstract describes offline learning of configuration trade-offs and online selection. Therefore, learning resource--performance trade-offs or transferring configuration knowledge is not the proposed novelty. Any adapted Scarf baseline must disclose differences in knobs, training effort, and available workloads.

Prediction and heterogeneous resource choice also have strong numerical precedents. MMP-SPS separates monitoring, predictor selection, and adaptation across multiple control windows, comparing prediction-error and bandit-based selection in Apache Storm~\cite{wladdimiro2026mmpsps}. Its evaluation also distinguishes resource savings from service degradation. Russo Russo et al. coordinate application- and operator-level controllers using Bayesian optimization and reinforcement learning for heterogeneous nodes~\cite{russorusso2023hierarchical}. Li et al.'s CEBFD approach combines resource-demand prediction with cost-aware executor placement for Spark Streaming~\cite{li2024cebfd}. The latter's micro-batch execution is not interchangeable with the proposed Flink setting; its relevance is the separation of demand prediction from placement cost.

ProOSStream combines offline scheduling optimization, federated reinforcement learning for online adaptation, and workload forecasting in edge--cloud environments~\cite{lajili2026proosstream}. It broadens the comparison beyond centralized scaling, but importing its distributed architecture into a single-host study would not establish equivalent deployment evidence. These works support testing whether diagnostic acquisition contributes beyond accurate prediction and competent numerical scheduling, with learning and controller overhead included.

\subsection{Partitioning, Scheduling, and Stateful Reconfiguration}
Performance deterioration can originate in data distribution, communication, or the mechanics of adaptation, not only insufficient CPU allocation. Ra-Stream combines resource scaling with communication-aware subgraph partitioning and fine-grained task scheduling~\cite{fan2026rastream}. SAP-DPO responds to fluctuating input with stream-aware partitioning, adaptive update cycles, and sliding-window parameter adjustment~\cite{ren2026sapdpo}. These approaches motivate separating capacity pressure from distribution-related problems when interpreting streaming telemetry. They also constrain a broad novelty claim about adaptive monitoring or update frequency.

CAOM detects bottleneck operators from runtime metrics and selects migration timing while considering interruption-related data accumulation~\cite{tan2025caom}. Thus, bottleneck-sensitive interventions and time-dependent migration costs already have direct streaming precedents. DRRS addresses the execution of stateful scaling through re-routing, record reordering, and subdivision of migrating state~\cite{qing2025drrs}. Its role differs from selecting a configuration: it reduces disruption when a scaling operation has already been chosen.

This distinction matters for EcoStream's transition model. Not every stateful adjustment should be assigned a mandatory full-restart cost; the model must reflect the actual execution path of the selected engine and configuration. Likewise, detecting skew or an impaired worker is not successful remediation unless the shared action catalogue contains an intervention that can address it. Cases outside that catalogue require an explicit no-change or escalation outcome, not an assumed successful repair.

Additional scheduling and recovery work makes this boundary more concrete. Ms-Stream coordinates grouping, partitioning, and placement to mitigate task- and node-level stragglers~\cite{wu2025msstream}. Ns-Stream instead prioritizes near-source processing through dynamic grouping, with its published implementation using Storm~\cite{wu2025nsstream}. Both inform the distinction between capacity and distribution bottlenecks; neither mechanism can be assumed available in an unmodified Flink executor. Incidents requiring unsupported redistribution should be excluded from a remediation claim or explicitly escalated.

Recovery scope is itself a design variable. Takdir et al. implement partial snapshots and local recovery in Flink to restore a subset of operators rather than the entire job~\cite{takdir2025localrecovery}. A-FP4S distributes state fragments over a peer-to-peer overlay and adapts parallel recovery to hardware heterogeneity~\cite{xu2023afp4s}. These mechanisms establish alternatives to global restoration. EcoStream would select among supported interventions, not claim to invent these recovery mechanisms; its measurements must use the actual restart, replay, and backlog costs observed in the deployed configuration.

\subsection{Agentic Diagnosis and Tool-Oriented Reasoning}
RCAgent is a particularly close precedent because autonomous operational diagnosis has already been applied to Alibaba Cloud Flink jobs~\cite{wang2024rcagent}. Its controller gathers information through tools, manages lengthy observations using snapshot keys, repairs malformed outputs, and can finalize a diagnosis. The existence of autonomous tool selection, efficient context management, or a stopping action is therefore not sufficient to distinguish EcoStream. The proposed difference concerns the criterion for acquiring evidence and stopping, and its effect on executed interventions.

ORACL couples language-model root-cause reasoning with resource-allocation recommendations~\cite{bai2026oracl}. Its inspected manuscript delegates numerical forecasting to a separate component and mediates actions through a verification interface. Diagnosis before allocation and separation of numerical prediction from language interpretation are consequently inherited design choices. EcoStream's planned comparison must ask whether another observation changes the preferred feasible action enough to justify its cost, rather than whether an agent can generate an allocation recommendation.

Scientific tool agents provide supporting execution patterns, but not evidence of streaming-control performance. SciToolAgent uses knowledge-graph-assisted retrieval for scientific tool selection and execution, including safety checks~\cite{ding2025scitoolagent}. BioMedAgent separates planning, coding, and execution, using interactive exploration and retrieval of successful experiences to construct biomedical workflows~\cite{bu2026biomedagent}. Their relevance is bounded to tool contracts, feedback, and experience organization. A large tool graph or additional agent roles should not be introduced into a small operational catalogue without evidence that the extra complexity improves decisions. Transfer from biomedical or scientific workflow success to latency-sensitive streaming control remains a hypothesis.

More recent diagnosis work strengthens the need for a precise comparison. TAMO's published scope explicitly combines tool-assisted LLM diagnosis with multimodal observations in cloud-native systems~\cite{zhang2025tamo}. LocaleXpert converts metrics, logs, and traces into language-based descriptions and coordinates specialist agents to localize and interpret failures~\cite{zhong2026localexpert}. MA-RCA separates diagnostic roles and introduces retrieval grounding and runtime hypothesis validation~\cite{fu2026marca}. Consequently, tool use, multimodal evidence, specialist agents, and hypothesis checking are established capabilities, not sufficient contribution statements on their own. Detailed reconciliation with TAMO's inaccessible final method remains necessary.

CausalRCA provides a non-LLM alternative, learning a weighted causal graph from metrics and ranking fine-grained candidate causes~\cite{xin2023causalrca}. Such diagnostic ranking does not by itself measure the benefit of an executed repair, and an observational causal model still requires its assumptions to be checked. The proposed study should separately measure diagnosis quality, action selection, investigation overhead, and post-intervention service recovery. It should not interpret a more convincing explanation as evidence of a better operational outcome.

\subsection{Sustainability and Measurement Boundaries}
Energy-aware scheduling is also established in Flink. Li et al. propose an SLA-based scheduling strategy together with coordinated inter-node and intra-node load balancing~\cite{li2025energyflink}. The reviewed abstract supports its scheduling motivation, but does not provide sufficient instrumentation detail for an independent physical-energy assessment. Resource allocation and energy consumption must therefore remain distinct quantities in comparisons.

Colangelo et al. demonstrate grid-responsive orchestration of AI workloads, using power controls, pausing, and resource reallocation while respecting designated performance-flexibility tiers~\cite{colangelo2026grid}. This provides a physical-measurement precedent, not a direct streaming-controller baseline: some AI-job service commitments permit throughput reductions, whereas EcoStream's proposed exact-query evaluation preserves semantics and offered load. Allocated core-seconds, consumed CPU, instantaneous power, total energy, and carbon emissions answer different questions. Without suitable physical instrumentation, the present study will target resource efficiency and service quality; sustainability remains motivation rather than a measured energy-saving claim.

\subsection{Comparative Gap and Proposed Study Question}
The reviewed systems already cover proactive adaptation, constrained optimization, selective profiling, autonomous diagnosis, and recovery- or migration-aware execution. Their strengths motivate a specific unresolved question for this study: under what conditions does acquiring an additional operational observation improve a stateful intervention enough to outweigh collection and waiting costs? This is a research question, not an assertion that every prior controller lacks decision-sensitive acquisition.

EcoStream would investigate an acquisition and stopping policy based on expected intervention loss. Its intended evaluation must separate this policy from the contributions of language interpretation and transition-cost modelling, using matched diagnostic tools and actions. A conventional diagnostic agent and deterministic processing of the same evidence are necessary controls. The proposed benefit would be weakened if cheap telemetry already determines the best action, if deterministic rules match the agent, or if investigation delays erase resource savings. These possible outcomes are part of the study, not exceptions to be removed after testing. Priority for the specific mechanism remains provisional until the closest final versions and implementations are reconciled.

\subsection{Scope of the Current Literature Evidence}
This is a targeted comparison, not an exhaustive systematic review. The additions comprise thirteen method papers and two surveys; eleven have 2024--2026 issue years, with four older journal contributions retained for directly relevant mechanisms. Some records were reviewed through publisher abstracts or author manuscripts; inaccessible method passages and differences from final published versions remain unresolved. TAMO is included using its published identity and title-level scope only, without importing uninspected preprint details. The bibliography identifies published versions, not ordinary preprints. Mechanism descriptions are limited to inspected support, and this draft neither reproduces literature-reported results nor presents them as EcoStream findings. The project's reference ledger records access and verification status; those pending checks prevent a submission-readiness claim.
